\section{Finite-field partitions and separated rows}
\label{sec:partitions}

We now construct partitions with few labels for which the functions of
Section~\ref{sec:compression} separate every pair of rows.  The construction
has two steps.  First, we choose directions so that, for each nonzero
symmetric multilinear form, one can make a form-dependent deletion of a
small fraction of the directions and obtain nonzero evaluations on every
tuple of remaining directions.  Second, these remaining directions
obstruct short paths in the corresponding partition graph.

We use the standard total-degree polynomial zero estimate
\cite[Corollary~1]{Schwartz1980}; related early work on randomized
polynomial evaluation is due to Zippel~\cite{Zippel1979}. We include
the precise finite-field statement and its elementary proof.

\begin{lemma}[Polynomial zero bound]
\label{lem:polynomial-zero}
Let $P\in\F_p[z_1,\ldots,z_k]$ be a nonzero polynomial of total degree at
most $d$, where $k,d$ are nonnegative integers and $p$ is prime.  If
$U_1,\ldots,U_k$ are independent and uniform in $\F_p$, then
\[
  \Pr\bigl(P(U_1,\ldots,U_k)=0\bigr)\le \frac{d}{p}.
\]
\end{lemma}

\begin{proof}
We induct on $k$.  When $k=0$, the polynomial is a nonzero constant,
and its zero probability is zero.  For $k\ge1$, write
\[
  P(z_1,\ldots,z_k)
  =\sum_{j=0}^{e}P_j(z_1,\ldots,z_{k-1})z_k^j,
  \qquad P_e\ne0.
\]
The total degree of $P_e$ is at most $d-e$, so the induction hypothesis
bounds the probability that this leading coefficient vanishes by
$(d-e)/p$.  Conditional on any values of the first $k-1$ variables for
which it does not vanish, the resulting univariate polynomial has degree
$e$ and at most $e$ roots, by successive division by linear factors.
Its conditional zero probability is therefore
at most $e/p$.  The sum of these two bounds is $d/p$.  This argument also
covers $e=0$ and, in particular, all nonzero constant polynomials.
\end{proof}

For positive integers $r,h$, set
\[
  D_j=\binom{r+j-1}{j}\quad(j\ge0).
\]
If $E=\F_p^r$, write $\operatorname{Sym}^j(E^*)$ for the space of
maps $E^j\to\F_p$ that are multilinear and invariant under permutations
of their arguments.  We interpret $\operatorname{Sym}^0(E^*)$ as
$\F_p$.  Relative to a basis $e_1,\ldots,e_r$, a symmetric $j$-linear
form is specified freely by its values on multisets of $j$ basis vectors.
The number of such multisets is $D_j$.  Thus
\begin{equation}
  \dim_{\F_p}\operatorname{Sym}^j(E^*)=D_j,
  \qquad
  \bigl|\operatorname{Sym}^j(E^*)\bigr|=p^{D_j}.
  \label{eq:symmetric-form-dimension}
\end{equation}

\begin{proposition}[Directions with nonvanishing evaluations]
\label{prop:robust-directions}
Let $0<\theta<1$ and let $r,h$ be positive integers.  Put
\begin{equation}
  w=2D_h,\qquad
  u=\left\lceil\frac{hw}{\theta}\right\rceil,
  \label{eq:field-parameters}
\end{equation}
and let $p$ be any prime satisfying
\begin{equation}
  p>\max\{h,h^2u^{2h}\}.
  \label{eq:field-prime-threshold}
\end{equation}
There exist $t_1,\ldots,t_u\in E=\F_p^r$ such that, for every nonzero
$Z\in\operatorname{Sym}^h(E^*)$, there is a set $T\subseteq[u]$ with
$|[u]\setminus T|\le\theta u$ for which
\begin{equation}
  Z(t_{i_1},\ldots,t_{i_h})\ne0
  \qquad\text{for all }(i_1,\ldots,i_h)\in T^h.
  \label{eq:robust-nonvanishing}
\end{equation}
Tuples in this assertion may have repeated indices.
\end{proposition}

\begin{proof}
Choose $t_1,\ldots,t_u$ independently and uniformly in $E$.  Fix a
nonzero form $Z$.  For $z=\sum_{a=1}^r z_a e_a$, its diagonal polynomial
is
\begin{equation}
  \begin{aligned}
  P_Z(z_1,\ldots,z_r)&=Z(z,\ldots,z)\\
  &=\sum_{\substack{\alpha_1+\cdots+\alpha_r=h\\\alpha_a\ge0}}
    \frac{h!}{\alpha_1!\cdots\alpha_r!}
    Z(e_1^{[\alpha_1]},\ldots,e_r^{[\alpha_r]})
    z_1^{\alpha_1}\cdots z_r^{\alpha_r},
  \end{aligned}
  \label{eq:diagonal-polynomial}
\end{equation}
where $e_a^{[\alpha_a]}$ means that $e_a$ is repeated $\alpha_a$ times
among the arguments.  At least one displayed value of $Z$ is nonzero.
Since $p>h$, every displayed multinomial coefficient is nonzero in
$\F_p$.  Distinct multisets give distinct monomials, so $P_Z$ is a
nonzero polynomial of total degree $h$.
This is the diagonal recovery of a symmetric multilinear form when
$h!$ is invertible; see \cite[Corollary~1.2]{FerreroMicali1979} and
\cite[Proposition~3.3]{DrapalVojtechovsky2009}. The coordinate argument
above gives the precise fact needed here.

For an ordered tuple $\mathbf i=(i_1,\ldots,i_h)\in[u]^h$, let
$\supp(\mathbf i)=\{i_1,\ldots,i_h\}$.  The expression
$Z(t_{i_1},\ldots,t_{i_h})$, considered as a polynomial in the scalar
coordinates of the vectors with indices in $\supp(\mathbf i)$, has total
degree $h$ and is nonzero: setting all those vector variables equal to
$z$ gives the nonzero polynomial \eqref{eq:diagonal-polynomial}.
Lemma~\ref{lem:polynomial-zero} therefore gives
\begin{equation}
  \Pr\bigl(Z(t_{i_1},\ldots,t_{i_h})=0\bigr)\le\frac hp.
  \label{eq:tuple-zero-probability}
\end{equation}
This applies to every pattern of repeated indices.

For this fixed $Z$, evaluations on tuples with pairwise disjoint index
supports depend on disjoint subsets of the independent vectors $t_i$.
Their zero events are consequently independent.  There are at most
$u^{hw}$ ordered lists of $w$ ordered $h$-tuples, and fewer than
$p^{D_h}$ choices of nonzero $Z$.  A union bound shows that the probability
that some nonzero $Z$ has $w$ zero tuples with pairwise disjoint supports
is at most
\begin{equation}
  p^{D_h}u^{hw}\left(\frac hp\right)^w
  =\left(\frac{h^2u^{2h}}p\right)^{D_h}<1.
  \label{eq:field-union-bound}
\end{equation}
The count includes all forms over the chosen field; no independence
between different forms is needed.  Fix vectors for which this event
does not occur.

For each nonzero $Z$, take a family of zero tuples with pairwise disjoint
supports that is maximal under inclusion.  It contains at most $w-1$
tuples.  Delete the union of their supports, and let $T$ be the remaining
indices.  The number deleted is at most
\[
  h(w-1)<hw\le\theta u.
\]
If any zero tuple were supported in $T$, its support would be disjoint
from all the selected supports, contradicting maximality.  Thus
\eqref{eq:robust-nonvanishing} holds.  Since $\theta<1$, the set $T$ is
nonempty.  The construction imposed no requirement that the sampled
vectors be distinct or nonzero.
\end{proof}

\begin{proposition}[Partitions with separated rows]
\label{prop:separated-rows}
With the parameters and directions of
Proposition~\ref{prop:robust-directions}, let
\[
  X=\operatorname{Sym}^h(E^*),\qquad
  L_i(x)=x(t_i,\cdot,\ldots,\cdot),
\]
and let $\mathcal T=\{T\subseteq[u]:|[u]\setminus T|\le\theta u\}$.
Then $|X|=m=p^{D_h}$, and each label map $L_i$ has at most
$q=p^{D_{h-1}}$ realized labels.  For the partition graphs of
Section~\ref{sec:compression}, every pair of distinct $x,y\in X$ has
some $T\in\mathcal T$ satisfying
\begin{equation}
  d_T(x,y)>h.
  \label{eq:graph-separation}
\end{equation}
Consequently, for $Y=X\times\mathcal T$ and
$R_x=(g_{(v,T)}(x))_{(v,T)\in Y}\in\R^Y$, the rows satisfy
\begin{equation}
  \|R_x-R_y\|_\infty=1\qquad(x\ne y).
  \label{eq:row-separation}
\end{equation}
\end{proposition}

\begin{proof}
The cardinality of $X$ follows from
\eqref{eq:symmetric-form-dimension}.  The contraction $L_i(x)$ is a
symmetric $(h-1)$-linear form, so its possible values belong to a space
of cardinality $p^{D_{h-1}}$.  This bounds the realized labels even if a
contraction is not surjective.

Fix distinct $x,y\in X$, and apply
Proposition~\ref{prop:robust-directions} to $Z=y-x$.  Let $T\in\mathcal T$
be the resulting nonempty set.  Recall that distinct vertices are
adjacent in the graph for $T$ if their labels agree at some $i\in T$.
Suppose there were a path
\[
  x=x_0,x_1,\ldots,x_k=y,\qquad 1\le k\le h.
\]
For each edge choose $i_j\in T$ such that
$L_{i_j}(x_j)=L_{i_j}(x_{j-1})$, and put
$\Delta_j=x_j-x_{j-1}$.  Thus
\[
  \Delta_j(t_{i_j},v_2,\ldots,v_h)=0
  \qquad\text{for every }v_2,\ldots,v_h\in E.
\]
Choose any $i_*\in T$, and form the $h$-tuple consisting of
$t_{i_1},\ldots,t_{i_k}$ followed by $h-k$ copies of $t_{i_*}$.
Every $\Delta_j$ vanishes on this tuple, since it is symmetric and one
of its arguments is $t_{i_j}$.  The identity
$Z=\sum_{j=1}^k\Delta_j$ implies that $Z$ also vanishes on the tuple,
contrary to \eqref{eq:robust-nonvanishing}.  This proves
\eqref{eq:graph-separation}; vertices in different components have
infinite distance and already satisfy it.

At the column $(x,T)$, the logarithmic profile gives
$g_{(x,T)}(x)=1$ and $g_{(x,T)}(y)=0$, because $d_T(x,y)>h$.
All column values lie in $[0,1]$, so this unit difference also gives
the exact equality \eqref{eq:row-separation}.
\end{proof}

The family $\mathcal T$ here contains every admissible subset and is
fixed with the partitions.  The separating set $T$ may depend on the
pair of rows.  This is compatible with Proposition~\ref{prop:compression},
whose common index may depend on the particular convex combination of
columns.  If $h$ also satisfies $\log3/\log(h+1)\le\theta$, the
compression bound applies with $q=p^{D_{h-1}}$ and $s=\lceil1/\theta\rceil$,
giving
\begin{equation}
  \log Q\le hsD_{h-1}\log p+C_{h,\theta,u}.
  \label{eq:field-compression-bound}
\end{equation}
Once $r,h,\theta$ are fixed, so are $u$ and $C_{h,\theta,u}$; the prime
$p$ can still be taken arbitrarily
large subject to \eqref{eq:field-prime-threshold}.
